% ==========================================================================
%  CODE_EXPLANATION_SHORT.tex  --  Condensed Walkthrough of the GA Prompt Optimiser
%  Target: ~32 pages.  Focus: call chains, function roles, cross-file data flow.
% ==========================================================================
\documentclass[11pt,a4paper]{report}

\usepackage[utf8]{inputenc}
\usepackage[T1]{fontenc}
\usepackage{lmodern}
\usepackage[margin=2.2cm,top=2.5cm,bottom=2.5cm]{geometry}
\usepackage{fancyhdr}
\usepackage{titlesec}
\usepackage{xcolor}
\usepackage{listings}
\usepackage{tcolorbox}
\tcbuselibrary{skins,breakable}
\usepackage{booktabs}
\usepackage{array}
\usepackage{enumitem}
\usepackage{hyperref}

% --- Colours ---
\definecolor{codebg}{HTML}{F7F7F7}
\definecolor{codeframe}{HTML}{CCCCCC}
\definecolor{kwcolor}{HTML}{0000CC}
\definecolor{strcolor}{HTML}{007700}
\definecolor{cmtcolor}{HTML}{888888}
\definecolor{conceptbg}{HTML}{EBF5FB}
\definecolor{conceptframe}{HTML}{2980B9}
\definecolor{warnbg}{HTML}{FFF3E0}
\definecolor{warnframe}{HTML}{E65100}
\definecolor{specbg}{HTML}{F0FFF0}
\definecolor{specframe}{HTML}{27AE60}
\definecolor{chainbg}{HTML}{FAFAFA}
\definecolor{chainframe}{HTML}{7F8C8D}

% --- Listings ---
\lstdefinestyle{py}{
  language=Python, basicstyle=\ttfamily\footnotesize,
  keywordstyle=\color{kwcolor}\bfseries, stringstyle=\color{strcolor},
  commentstyle=\color{cmtcolor}\itshape,
  backgroundcolor=\color{codebg}, frame=single, rulecolor=\color{codeframe},
  breaklines=true, showstringspaces=false,
  numbers=left, numberstyle=\tiny\color{gray}, numbersep=6pt,
  tabsize=4, xleftmargin=1.2em, framexleftmargin=1.2em,
  aboveskip=0.5em, belowskip=0.5em,
  literate={---}{{---}}1,
}
\lstdefinestyle{chain}{
  language={}, basicstyle=\ttfamily\footnotesize,
  backgroundcolor=\color{chainbg}, frame=single, rulecolor=\color{chainframe},
  breaklines=true, showstringspaces=false, numbers=none,
  aboveskip=0.4em, belowskip=0.4em,
  xleftmargin=0.3em, framexleftmargin=0.3em,
}
\lstset{style=py}

% --- Boxes ---
\newtcolorbox{concept}[1]{colback=conceptbg,colframe=conceptframe,
  fonttitle=\bfseries,title={#1},boxrule=0.7pt,arc=3pt,
  left=5pt,right=5pt,top=3pt,bottom=3pt,breakable}
\newtcolorbox{warn}[1]{colback=warnbg,colframe=warnframe,
  fonttitle=\bfseries,title={#1},boxrule=0.7pt,arc=3pt,
  left=5pt,right=5pt,top=3pt,bottom=3pt,breakable}
\newtcolorbox{spec}[1]{colback=specbg,colframe=specframe,
  fonttitle=\bfseries\ttfamily,title={#1},boxrule=0.7pt,arc=3pt,
  left=5pt,right=5pt,top=3pt,bottom=3pt,breakable}

% --- Macros ---
\newcommand{\fn}[1]{\texttt{#1}}
\newcommand{\pyf}[1]{\texttt{#1}}

% --- Headers ---
\pagestyle{fancy}\fancyhf{}
\fancyhead[L]{\small\leftmark}
\fancyhead[R]{\small Code Explanation (Condensed)}
\fancyfoot[C]{\thepage}
\renewcommand{\headrulewidth}{0.4pt}

% --- Chapter style ---
\titleformat{\chapter}[display]
  {\normalfont\huge\bfseries}{\chaptertitlename\ \thechapter}{14pt}{\Huge}
\titlespacing*{\chapter}{0pt}{-20pt}{20pt}

% ==========================================================================
\begin{document}

% --- Title page ---
\begin{titlepage}
\centering\vspace*{2.5cm}
{\Huge\bfseries Code Explanation\par}\vspace{0.6cm}
{\LARGE GA Prompt Optimiser for LLM Code Generation\par}\vspace{1cm}
{\Large Evolving System Prompts on Codeforces Problems\par}\vspace{2cm}
{\large Group 1 --- Computational Intelligence, Spring 2026\par}\vspace{0.8cm}
{\large April 2026\par}\vfill
{\small Condensed walkthrough: every file, every function, every call chain.\\
Backend: Google Cloud Vertex AI (Gemini 2.5 Flash \& Pro).\\
Includes cross-file data-flow diagrams and complete worked examples.}
\end{titlepage}

\tableofcontents\newpage

% ==========================================================================
\chapter{Introduction \& Architecture}
% ==========================================================================

\section{What the Project Does}

A \textbf{Genetic Algorithm (GA)} automatically discovers the best \textbf{system
prompt} --- a paragraph of English instructions --- for a Gemini model that generates
\textbf{C++ code}.  Each prompt is a \emph{chromosome}; fitness is the average fraction
of test cases the LLM-generated C++ program passes on 65 Codeforces coding problems.
Over 12 generations of selection, crossover, and mutation the GA evolves a prompt that
solves more problems correctly.

\begin{concept}{Key idea}
Manual prompt engineering is trial-and-error.  This project automates it: the GA
\emph{searches} all possible English instruction paragraphs, using test-case pass
rates as the fitness signal.  The LLM never sees the test cases --- only the problem
statement.
\end{concept}

\section{The Nine Files}

\begin{center}
\begin{tabular}{@{}lrp{7cm}@{}}
\toprule
\textbf{File} & \textbf{Lines} & \textbf{Role} \\
\midrule
\pyf{config.py}        & 83  & Single source of truth for all hyperparameters and paths \\
\pyf{seeds.py}         & 137 & 15 hand-crafted generation-0 chromosomes (C++ focused) \\
\pyf{dataset.py}       & 253 & Load 305 Codeforces problems; stratified split into evolution / holdout \\
\pyf{evaluator.py}     & 295 & Extract C++ code; compile with g++ -std=c++23; run tests \\
\pyf{cost\_tracker.py} & 244 & Stream token counts to JSONL; compute Vertex AI cost; write JSON reports \\
\pyf{llm.py}           & 470 & All Vertex AI Gemini calls (code gen, crossover, mutation) \\
\pyf{ga.py}            & 717 & The evolutionary loop: select, breed, evaluate (parallel), replace; 3 selection schemes \\
\pyf{main.py}          & 545 & Orchestrate the run; checkpoint; resume; holdout; genealogy \\
\pyf{genealogy.py}     & 154 & Trace ancestry of best evolved prompt through checkpoint files \\
\bottomrule
\end{tabular}
\end{center}

\section{Module Dependency Graph}

\begin{lstlisting}[style=chain]
main.py ----imports----> config, dataset, evaluator, ga, llm
                         cost_tracker (singleton: tracker)
                         seeds        (SEED_PROMPTS list)
                         genealogy    (generate_report)

ga.py   ----imports----> config, evaluator, llm

llm.py  ----imports----> config
                         cost_tracker (singleton: tracker)
                         google.genai (Vertex AI SDK)

evaluator.py   -------> stdlib only  (os, re, signal, subprocess, tempfile, time)
dataset.py     -------> stdlib only  (json, random, sys, pathlib)
cost_tracker.py ------> stdlib only  (json, threading, time, dataclasses, datetime, pathlib)
genealogy.py   -------> stdlib only  (json, pathlib, collections)
seeds.py       -------> no imports
\end{lstlisting}

\begin{concept}{Network boundary}
\pyf{llm.py} is the \textbf{only} file that touches the network.  It uses the
\fn{google-genai} Python SDK with Google Cloud Vertex AI backend, authenticated via a
service account JSON file (\fn{service-account.json}).  No interactive login or API key
is required --- the service account credential is loaded once at startup.
All other files are pure computation or local disk I/O.
\end{concept}

\section{Two-Model Architecture}

\begin{center}
\begin{tabular}{@{}llll@{}}
\toprule
\textbf{Role} & \textbf{Model} & \textbf{Temp.} & \textbf{Purpose} \\
\midrule
Target (subject)    & \fn{gemini-2.5-flash} & 0.0 (enforced) & Generate C++ code; fitness eval \\
Optimizer (tool)    & \fn{gemini-2.5-pro}   & 0.7 (enforced) & Crossover and mutation operators \\
\bottomrule
\end{tabular}
\end{center}

\begin{concept}{Temperature enforcement}
Unlike the previous Claude CLI design (where temperature flags were unavailable),
the Vertex AI SDK enforces both temperatures exactly via
\fn{GenerateContentConfig(temperature=...)}.  \fn{TARGET\_TEMPERATURE=0.0} makes
code generation deterministic; \fn{OPTIMIZER\_TEMPERATURE=0.7} injects creative
variation into the genetic operators.
\end{concept}

\begin{concept}{Thinking tokens}
Gemini 2.5 models generate internal chain-of-thought (\emph{thinking tokens}) before
producing the visible response.  These are billed at the same output rate and are
tracked separately via \fn{usage\_metadata.thoughts\_token\_count}.  Budgets:
\fn{thinking\_budget=512} for fitness calls (Flash), \fn{thinking\_budget=2048} for
crossover/mutation (Pro).
\end{concept}

\section{Master Call-Flow (one generation)}

\begin{lstlisting}[style=chain]
main.main()
  |-- Vertex AI preflight: verify service-account.json exists + client connects
  |-- dataset.load_problems()     -> list[dict]  (305 Codeforces problems)
  |-- dataset.split(problems, 65, 200, 42)       -> (evolution[65], holdout[200])
  |-- tracker.set_log_path(run_dir/"calls.jsonl")
  |-- tracker.set_live_path(run_dir/"costs_live.json")
  |-- tracker.set_costs_live_path(COSTS_DIR/"cost_<run_dir_name>.json")
  |-- llm.set_error_log_path(run_dir/"llm_errors.log")
  |
  +-- ga.run(SEED_PROMPTS[:5], evolution_pool, run_dir, on_generation=save_checkpoint,
             selection_scheme=args.scheme, elitism_count=1 if scheme=="elitism" else None)
        |
        | [GEN 0] for each of 5 seeds (sequential):
        |   save seed_NN_<uid>.json checkpoint first (crash recovery)
        |   compute_fitness(gen=0, ind, problems)   [65 problems, PARALLEL_EVALS=5]
        |     ThreadPoolExecutor(max_workers=5):
        |       _eval_one_problem(gen, ind, prob, idx, total)
        |         time.sleep(random.uniform(0,1))   [stagger concurrent requests]
        |         -> llm.generate_code(ind.prompt, prob["question"])
        |              client.models.generate_content(...)  [Gemini Flash]
        |              -> tracker.record(phase, model, tokens, cost)
        |         -> _eval_code_for_problem(gen, ind, prob, idx, total, response, llm_time)
        |              evaluator.extract_code(response)  -> C++ code string
        |              evaluator.run_tests_verbose(code, inputs, outputs, timeout=2)
        |                _compile(code)          -> g++ -O2 -std=c++23 [Popen+setsid]
        |                _run_one(bin, stdin, 2) -> actual stdout       [Popen+setsid]
        |                _outputs_match(actual, expected)
        |              _log_evaluation(entry)    -> evaluations_genNN.jsonl
        |              returns (score, failure_example_or_None)
        |       _assign_failure_examples(ind, collected_failures) -> ind.failure_examples
        |     fitness = mean(scores)
        |
        | [GENS 1..12] for each generation:
        |   _select(population, rng)  x2 -> parents  [tournament/roulette per --scheme]
        |   breed(p1, p2, rng, allow_mutation=True or False)  -> offspring
        |     llm.crossover(a, b)                    [80% chance]
        |     llm.mutate_inject(p, failure_examples) [25%] (if allow_mutation)
        |     llm.mutate_delete(p)                   [15%] (if allow_mutation)
        |     llm.mutate_rephrase(p)                 [10%] (if allow_mutation)
        |   Note: first CROSSOVER_ONLY_SLOTS=2 offspring skip mutations
        |   compute_fitness(gen, offspring, problems) [parallel, 65 problems each]
        |   save_checkpoint(run_dir, gen, pop, stats) -> generation_NN.json
        |
        | Convergence check each gen:
        |   CV = std/mean; if CV < 0.01 for 3 consecutive gens -> early stop
        |
        +-- returns (best_Individual, history[list of dicts])
  |
  |-- evaluate_on_holdout(best.prompt, holdout[200])  [ThreadPoolExecutor, parallel]
  |-- genealogy.generate_report(run_dir, best.uid)    -> genealogy_report.txt
  +-- tracker.flush(extra)   -> Code/Costs/cost_TIMESTAMP.json
\end{lstlisting}

\begin{concept}{API call scale}
One full run (no early stop) makes approximately:\\
$5 \times 65 = 325$ Flash calls (gen 0) +
$12 \times 5 \times 65 = 3{,}900$ Flash calls (gens 1--12) +
$200$ Flash calls (holdout) = \textbf{4,425 Gemini Flash calls}.\\
Crossover/mutation: $\approx 12 \times (5 \times 0.8) + (3 \times 0.25) + (3 \times 0.15) + (3 \times 0.10) = 66$ Pro calls.\\
Estimated cost: $\approx\$2.12$ (Flash) $+$ $\approx\$0.89$ (Pro) $\approx \$3.01$ total.
\end{concept}


% ==========================================================================
\chapter{\pyf{config.py} --- Hyperparameters}
% ==========================================================================

\section{Purpose \& Design}

Every number and path lives here.  No other file contains hard-coded values.
All constants are read at import time; changing \pyf{config.py} affects the whole system.
Authentication uses a GCP service account JSON file --- no API key or \texttt{.env}
file is required.

\section{Key Constants (current values)}

\begin{lstlisting}[style=py, caption={config.py (complete)}]
from pathlib import Path

ROOT         = Path(__file__).parent
APPS_TRAIN   = ROOT / "Dataset" / "APPS" / "train"
APPS_TEST    = ROOT / "Dataset" / "APPS" / "test"
RESULTS_DIR  = Path(__file__).parent / "results"

# Dataset -- 305 pre-selected Codeforces problems
EVOLUTION_SIZE      = 65    # stratified evolution pool (5 harder + 60 regular)
HOLDOUT_SIZE        = 200   # held-out for final generalisation test
RANDOM_SEED         = 42
EARLY_STOP_FAILURES = 3     # stop running tests after N consecutive failures

# Genetic Algorithm
POPULATION_SIZE      = 5
GENERATIONS          = 12
ELITISM_COUNT        = 0    # no frozen elites -- all slots bred via crossover+mutation
CROSSOVER_ONLY_SLOTS = 2    # first N offspring per gen get crossover only (no mutation)
TOURNAMENT_SIZE      = 3    # best-of-3 random draws for each parent

CROSSOVER_RATE    = 0.8
MUT_INJECT_PROB   = 0.25
MUT_DELETE_PROB   = 0.15
MUT_REPHRASE_PROB = 0.10

# Google Cloud Vertex AI -- Gemini backend
SERVICE_ACCOUNT_PATH = Path(__file__).parent / "service-account.json"
GCP_REGION           = "us-central1"

TARGET_MODEL    = "gemini-2.5-flash"    # code generation (fitness eval)
OPTIMIZER_MODEL = "gemini-2.5-pro"      # crossover / mutation meta-prompts

TARGET_TEMPERATURE    = 0.0   # enforced via GenerateContentConfig
OPTIMIZER_TEMPERATURE = 0.7   # enforced via GenerateContentConfig

API_FALLBACK_MAX_TOKENS_CODE = 8192   # max output tokens for code generation
API_FALLBACK_MAX_TOKENS_OPS  = 700    # crossover / mutation outputs are short

# Convergence detection (CV = std/mean per generation)
CONVERGENCE_CV_THRESHOLD = 0.01   # 1% variation -> converged
CONVERGENCE_PATIENCE     = 3      # consecutive stagnant generations to confirm
CONVERGENCE_MIN_MEAN     = 0.01   # only check CV when mean fitness > this

# Code execution (g++ compilation + binary runs)
EXEC_TIMEOUT    = 2    # seconds per test-case binary run
COMPILE_TIMEOUT = 15   # seconds allowed for g++ compilation
PARALLEL_EVALS  = 5    # problems evaluated simultaneously per individual
\end{lstlisting}

\begin{concept}{Path change: ROOT is now \fn{Path(\_\_file\_\_).parent}}
\fn{ROOT} previously resolved to the project parent directory (\fn{parent.parent}).
It now resolves to the \fn{Code/} directory itself (\fn{parent}).
Consequently the dataset path changed from \fn{ROOT / "Dataset Examples" / ...}
to \fn{ROOT / "Dataset" / ...} to match the actual folder layout.
This means all dataset files live under \fn{Code/Dataset/APPS/\{train,test\}/}.
\end{concept}

\begin{concept}{CROSSOVER\_ONLY\_SLOTS and elitism change}
\fn{ELITISM\_COUNT} was lowered from 2 to 0 --- no individuals are frozen across
generations.  Instead, \fn{CROSSOVER\_ONLY\_SLOTS=2} ensures the first 2 offspring per
generation undergo only crossover (no mutation), providing stability without hard-locking
elite individuals.  The remaining 3 offspring per generation are eligible for all mutation
operators.
\end{concept}

\begin{concept}{Mutations are independent}
Each of the three mutation probabilities is checked with a separate \fn{rng.random()}
call, so all three can fire on the same offspring (probability $0.25 \times 0.15 \times
0.10 = 0.00375$), or none can fire ($0.75 \times 0.85 \times 0.90 = 0.574$).
Only offspring beyond \fn{CROSSOVER\_ONLY\_SLOTS} are eligible for mutation.
\end{concept}


% ==========================================================================
\chapter{\pyf{seeds.py} --- Initial Population}
% ==========================================================================

\pyf{seeds.py} defines \fn{SEED\_PROMPTS}: a \texttt{list[str]} of 15 diverse
instruction paragraphs.  All prompts require a \textbf{complete, compilable C++ program}
as output; each emphasises a different cognitive strategy:

\begin{center}
\begin{tabular}{@{}rll@{}}
\toprule
\# & Strategy & Core emphasis \\
\midrule
1  & Structured analysis     & Analyse input/output format, constraints, edge cases before coding \\
2  & Terse and direct        & Full C++ solution; compile with \fn{g++ -std=c++23}; just working code \\
3  & Competitive programmer  & Identify algorithm pattern (greedy, DP, graph, binary search, math) \\
4  & Step-by-step reasoning  & 4-step: understand, identify algorithm, handle boundary, implement \\
5  & Example-driven          & Trace sample inputs manually; confirm behaviour before implementing \\
6  & Conversational          & Give me the full C++ source with main; handle all edge cases \\
7  & Edge-case defender      & n=0, n=1, all equal, maximum values, integer overflow \\
8  & Algorithm-first         & Determine algorithm + time/space complexity before writing code \\
9  & Specification-heavy     & 5-point spec: cin/cout, \fn{g++ -O2 -std=c++23}, all cases, exact format \\
10 & Pseudocode-then-code    & English sketch first, then translate to C++; cin/cout; include main() \\
11 & Constraint-aware        & If n=$10^5$, $O(n^2)$ will TLE; match complexity to constraints \\
12 & Defensive / robust      & Assume adversarial inputs at extreme constraint boundaries \\
13 & Readability-first       & Meaningful variable names; correct beats clever; standard I/O \\
14 & Verify-against-examples & Mentally verify against sample inputs before writing code \\
15 & Direct and focused      & Entire C++ program; reads stdin; writes stdout; passes all tests \\
\bottomrule
\end{tabular}
\end{center}

The file ends with \fn{assert len(SEED\_PROMPTS) >= 5} (previously \fn{== 15}).
The guard now allows populations smaller than 15 --- \fn{main.py} uses only
\fn{SEED\_PROMPTS[:POPULATION\_SIZE]} (the first 5 with \fn{POPULATION\_SIZE=5}).
Seeds 2 and 9 reference \fn{g++ -std=c++23} (updated from \fn{c++17}).

In \fn{ga.run()}, the selected seeds become \fn{Individual} objects:
\begin{lstlisting}[style=py]
population = [Individual(prompt=p, operators="seed")
              for p in seed_prompts[:config.POPULATION_SIZE]]
# Each Individual: prompt=<seed>, fitness=-1.0, uid=<random 8-char hex>,
#                  parent_a="", parent_b="", operators="seed",
#                  intermediate_steps={}, failure_examples=[]
\end{lstlisting}


% ==========================================================================
\chapter{\pyf{dataset.py} --- Loading \& Splitting Problems}
% ==========================================================================

\section{Overview}

\pyf{dataset.py} loads 305 pre-selected Codeforces problems (stdin/stdout, C++ programs)
and performs a \textbf{stratified split} into an evolution pool (65) and a held-out test
set (200).  Problem IDs are hard-coded in module-level dictionaries, ensuring
reproducibility without requiring a specific dataset directory structure.

\begin{concept}{Why stratified?}
A plain shuffle risks all harder problems landing in holdout.  Stratified sampling
draws a fixed count from each difficulty tier, guaranteeing the evolution pool sees
problems at every rating level --- including the harder 1800/1900-rated problems
added in this version.
\end{concept}

\section{\fn{load\_problems()} $\to$ \fn{list[dict]}}

\begin{spec}{load\_problems() -> list[dict]}
Called by: main.main() \\
Calls: Path.exists(), json.loads(), Path.read\_text()
\end{spec}

Iterates over the hard-coded \fn{\_SELECTED} dict (split $\to$ rating $\to$ folder IDs).
For each problem folder it reads \fn{question.txt} and \fn{input\_output.json}.
Skips problems with missing files, bad JSON, or empty test cases.

Rating distribution (305 problems total):

\begin{center}
\begin{tabular}{@{}lrrrr@{}}
\toprule
\textbf{Tier} & \textbf{Rating} & \textbf{Train} & \textbf{Test} & \textbf{Total} \\
\midrule
Regular & 800  & 29 & 21 & 50 \\
Regular & 1000 & 10 & 40 & 50 \\
Regular & 1100 & 14 & 19 & 33 \\
Regular & 1200 & 23 & 10 & 33 \\
Regular & 1300 & 23 & 11 & 34 \\
Regular & 1400 & 27 &  5 & 32 \\
Regular & 1500 & 33 &  0 & 33 \\
Regular & 1600 & 17 &  0 & 17 \\
Regular & 1700 & 18 &  0 & 18 \\
Harder  & 1800 &  3 &  0 &  3 \\
Harder  & 1900 &  2 &  0 &  2 \\
\midrule
\textbf{Total} & & 199 & 106 & 305 \\
\bottomrule
\end{tabular}
\end{center}

Each passing problem produces:
\begin{lstlisting}[style=py]
{
    "id":       "train/0002",   # "<split>/<folder>"
    "question": "Given n integers, ...",
    "inputs":   ["3\n1 2 3\n", "1\n0\n"],   # raw stdin strings
    "outputs":  ["6\n",        "0\n"],        # expected stdout strings
    "rating":   800,
}
\end{lstlisting}

\section{\fn{split()} $\to$ \fn{(evolution, holdout, reserved)}}

\begin{spec}{split(problems, evolution\_size=65, holdout\_size=200, seed=42) -> (list, list, list)}
Called by: main.main() \\
Calls: random.Random(seed), list comprehension, rng.shuffle(), rng.sample()
\end{spec}

The stratified split algorithm:

\begin{lstlisting}[style=chain]
split(problems, evolution_size=65, holdout_size=200, seed=42):

  rng = random.Random(seed)   # isolated, reproducible
  problems_by_rating = group problems by their "rating" field

  [HARDER PROBLEMS -- always go to evolution]
  harder_tiers = {1800, 1900}
  harder_problems = [p for p in problems if p["rating"] in harder_tiers]
  evolution = list(harder_problems)      # 5 problems (3+2)

  [REGULAR TIERS -- stratified draw]
  _EVOLUTION_PER_TIER = {800:7, 1000:6, 1100:6, 1200:7, 1300:7,
                          1400:7, 1500:7, 1600:6, 1700:7}   # sums to 60
  for each regular rating tier:
      tier_problems = [p for p in problems if p["rating"] == rating]
      rng.shuffle(tier_problems)    # seeded shuffle
      evolution += tier_problems[:count]   # take exactly count from this tier

  # Result: evolution has exactly 65 problems (5 harder + 60 regular)

  [HOLDOUT -- remaining problems, shuffled]
  remaining = [p for p in problems if p not in evolution set]
  rng.shuffle(remaining)
  holdout  = remaining[:holdout_size]      # 200 problems
  reserved = remaining[holdout_size:]      # leftovers (small number)
  return (evolution, holdout, reserved)
\end{lstlisting}

\section{\fn{reserved\_train\_problems(problems, seed=42)} $\to$ \fn{list[dict]}}

\begin{spec}{reserved\_train\_problems(problems, seed=42) -> list[dict]}
Called by: main.main() (optional, for analysis) \\
Returns: the 35 train-split problems not selected for evolution
\end{spec}

This function returns the train-split problems that were shuffled out of the
65-problem evolution pool --- useful for post-hoc analysis without touching the
holdout set.

\begin{concept}{Fixed seed is critical}
The same seed 42 must be used for all \fn{split()} calls throughout a run.  If the
65 evolution problems changed between generations, fitness scores from generation 3
and generation 7 would not be comparable.  The seed is stored in \fn{config.json}
in each run directory.
\end{concept}


% ==========================================================================
\chapter{\pyf{evaluator.py} --- C++ Extraction, Compilation \& Testing}
% ==========================================================================

\section{Overview}

Two responsibilities: (1) strip C++ code out of an LLM's raw text response, and
(2) compile it with \fn{g++} and run the resulting binary against test cases, returning
a partial-credit score.  Uses only the Python standard library, plus \fn{signal} for
process-group killing.

\section{\fn{extract\_code(response: str)} $\to$ \fn{str}}

\begin{spec}{extract\_code(response) -> str}
Called by: ga.\_eval\_code\_for\_problem(); main.evaluate\_on\_holdout()
\end{spec}

Tries three strategies in order:

\begin{lstlisting}[style=py]
# 1. Tagged C++ fence:  ```cpp ... ``` or ```c++ ... ```
m = re.search(r"```(?:cpp|c\+\+)\s*\n([\s\S]*?)```", response, re.IGNORECASE)
if m: return m.group(1).strip()

# 2. Any fence:  ``` ... ```
m = re.search(r"```\s*\n?([\s\S]*?)```", response)
if m: return m.group(1).strip()

# 3. Assume raw code (no fences)
return response.strip()
\end{lstlisting}

\section{Process-Group Management: \fn{\_kill\_group()}}

\begin{spec}{\_kill\_group(proc) -> None}
Called by: \_compile(), \_run\_one() on timeout
\end{spec}

\begin{lstlisting}[style=py]
def _kill_group(proc):
    try:
        os.killpg(os.getpgid(proc.pid), signal.SIGKILL)
    except ProcessLookupError:
        pass
\end{lstlisting}

\fn{SIGKILL} is sent to the entire \textbf{process group}, not just the parent.
This prevents orphaned child processes (e.g., a spawned shell or forked worker)
from persisting after a timeout.  Both \fn{\_compile()} and \fn{\_run\_one()}
create new process groups via \fn{preexec\_fn=os.setsid} in \fn{subprocess.Popen}.

\section{\fn{\_compile(code: str)} $\to$ \fn{(bin\_path | None, error\_msg)}}

\begin{spec}{\_compile(code) -> (str | None, str)}
Called by: run\_tests\_verbose() \\
Calls: tempfile.mkstemp(), subprocess.Popen, proc.communicate(timeout=COMPILE\_TIMEOUT)
\end{spec}

\begin{lstlisting}[style=chain]
_compile(code):
  src_fd, src_path = tempfile.mkstemp(suffix=".cpp")   # temp .cpp file
  bin_path = src_path[:-4]                              # strip extension

  write code to src_path
  proc = subprocess.Popen(
      ["g++", "-O2", "-std=c++23", "-o", bin_path, src_path],
      stdout=subprocess.PIPE, stderr=subprocess.PIPE,
      preexec_fn=os.setsid)                # new process group for kill_group()

  try:
      _, stderr = proc.communicate(timeout=COMPILE_TIMEOUT)
  except subprocess.TimeoutExpired:
      _kill_group(proc)
      proc.communicate()
      return None, "compile_timeout"

  if proc.returncode != 0:
      return None, stderr.decode()[:500]  # compile error
  return bin_path, ""                     # success

  finally: os.unlink(src_path)            # always delete .cpp
           if failed: os.unlink(bin_path) # only delete binary on failure
\end{lstlisting}

\begin{warn}{g++ -std=c++23 (not c++17)}
The compiler standard was updated from \fn{-std=c++17} to \fn{-std=c++23}.
Seeds 2 and 9 in \fn{seeds.py} were correspondingly updated to reference \fn{c++23}.
Modern Codeforces problems increasingly use C++23 features.
\end{warn}

\section{\fn{run\_tests\_verbose()} $\to$ \fn{dict}}

\begin{spec}{run\_tests\_verbose(code, inputs, outputs, timeout=2) -> dict}
Called by: ga.\_eval\_code\_for\_problem() \\
Calls: \_compile(), \_run\_one(), \_outputs\_match()
\end{spec}

\begin{lstlisting}[style=chain]
run_tests_verbose(code, inputs, outputs, timeout):

  if not code or not code.strip():
      return {score:0.0, reason:"empty_code", ...}

  [COMPILE]
  bin_path, err = _compile(code)
  if bin_path is None:
      return {score:0.0, reason:"compile_fail", compile:{error:err}, ...}

  [RUN TESTS]
  passed = 0; consecutive_fails = 0

  for i in range(len(inputs)):
      actual = _run_one(bin_path, inputs[i], timeout)
      is_pass = actual is not None and _outputs_match(actual, outputs[i])

      record test_cases[i] = {index, passed, actual, expected, time_s, timeout}

      if is_pass: passed += 1; consecutive_fails = 0
      else:
          consecutive_fails += 1
          if consecutive_fails >= EARLY_STOP_FAILURES:
              early_stopped = True; break

  score = passed / len(inputs)    # denominator = TOTAL tests (fair penalty)
  reason = "perfect" | "partial_pass" | "all_fail" | "all_fail_early_stop"

  return {score, reason, compile, test_cases, tests_total, tests_run,
          early_stopped, total_test_time_s}
\end{lstlisting}

\fn{run\_tests()} is a thin wrapper returning only the float score:
\begin{lstlisting}[style=py]
def run_tests(code, inputs, outputs, timeout=2) -> float:
    return run_tests_verbose(code, inputs, outputs, timeout)["score"]
\end{lstlisting}

\begin{concept}{Early stopping}
After \fn{EARLY\_STOP\_FAILURES=3} consecutive failures, the test loop exits.  The
score denominator is still \fn{len(inputs)} (total), so skipped tests count as
failures.  This saves wall-clock time on clearly broken solutions.
\end{concept}


% ==========================================================================
\chapter{\pyf{cost\_tracker.py} --- Token \& Cost Logging}
% ==========================================================================

\section{Purpose \& Design}

Logs token counts and estimated costs for every Vertex AI call.  The tracker streams
each call record to a JSONL file immediately (no in-memory list), keeping memory
usage flat regardless of run length.  Three live JSON files are updated after every
call; a final summary is written to \fn{Code/Costs/} at run end.

\begin{concept}{Vertex AI billing}
The project uses Google Cloud Vertex AI.  Costs are real GCP charges --- no flat
subscription.  The \fn{PRICING} table contains verified Gemini rates (April 2026).
Gemini 2.5 Flash: \$0.15/\$0.60 per MTok (in/out).
Gemini 2.5 Pro: \$1.25/\$10.00 per MTok (in/out).
The \$300 GCP free credit applies.
\end{concept}

\section{\fn{CallRecord} --- One Call's Data}

\begin{lstlisting}[style=py]
PRICING = {   # (input_$/MTok, output_$/MTok)
    "gemini-2.5-flash":     (0.15,  0.60),
    "gemini-2.5-pro":       (1.25, 10.00),
    "gemini-2.0-flash":     (0.10,  0.40),
    "gemini-2.0-flash-001": (0.10,  0.40),
    "claude-haiku-4-5@20251001":  (1.00,  5.00),  # kept for reference
    "claude-haiku-4-5-20251001":  (1.00,  5.00),
    "claude-sonnet-4-6":          (3.00, 15.00),
}

@dataclass
class CallRecord:
    phase: str; model: str
    input_tokens: int; output_tokens: int
    cost_usd: float | None = None   # when set, used directly

    @property
    def total_cost(self) -> float:
        if self.cost_usd is not None:
            return self.cost_usd
        rate_in, rate_out = PRICING.get(self.model, (0, 0))
        return (self.input_tokens * rate_in
              + self.output_tokens * rate_out) / 1_000_000
\end{lstlisting}

\section{\fn{CostTracker} --- The Streaming Singleton}

\begin{spec}{class CostTracker}
Module-level singleton: tracker = CostTracker() \\
Imported by llm.py as: from cost\_tracker import tracker
\end{spec}

\begin{lstlisting}[style=chain]
CostTracker.set_log_path(path)
  Called by: main.main() before GA starts
  Action: opens JSONL stream to calls.jsonl in the run directory

CostTracker.set_live_path(path)
  Called by: main.main()
  Action: path updated with current totals after every call (run_dir/costs_live.json)

CostTracker.set_costs_live_path(path)
  Called by: main.main()
  Action: second live path in Costs/ folder (Costs/run_live.json) updated every call

CostTracker.restore_from_jsonl(path) -> int
  Called by: main.main() on resume (--resume flag)
  Action: replays existing calls.jsonl to restore running totals; returns lines loaded

CostTracker.record(phase, model, input_tokens, output_tokens, cost_usd, extra)
  Called by: llm.py after every successful Vertex AI response
  Thread-safe (internal lock)
  Action:
    - updates running totals (_total_calls, _total_cost_usd, etc.)
    - writes one JSON line immediately to calls.jsonl (no in-memory list)
    {"t":"HH:MM:SS", "phase":..., "model":..., "in":..., "out":...,
     "cost":..., "total_calls":..., "running_cost":...}
    - overwrites both live JSON files

CostTracker.flush(extra=None) -> Path
  Called by: main.main() at the very end
  Action:
    - assembles summary from running totals (_phase_totals, _model_totals)
    - writes Code/Costs/cost_TIMESTAMP.json
    - returns Path of the written file
\end{lstlisting}

\begin{concept}{Streaming vs.\ in-memory}
Each record is written to disk immediately and discarded; running totals are maintained
as simple scalars.  A 4,425-call run uses negligible memory for cost tracking.
The \fn{extra} parameter to \fn{record()} lets crossover/mutation calls embed
input/output prompt text in the JSONL for audit purposes.
\end{concept}


% ==========================================================================
\chapter{\pyf{llm.py} --- All Vertex AI Calls}
% ==========================================================================

\section{Overview}

The sole interface to Gemini.  All public functions route through the private
\fn{\_call\_realtime\_api()} wrapper, which uses the \fn{google-genai} SDK to call
Vertex AI, handles exponential-backoff retries, tracks thinking tokens, and logs costs.

\begin{lstlisting}[style=chain]
Public API:
  generate_code(system_prompt, question)              -> str  [C++ code text]
  crossover(prompt_a, prompt_b)                       -> str  [merged prompt]
  mutate_inject(prompt, failure_examples=None)        -> str  [prompt + new strategy]
  mutate_delete(prompt)                               -> str  [prompt - weakest sentence]
  mutate_rephrase(prompt)                             -> str  [rephrased prompt]
  set_error_log_path(path)                            -> None [llm_errors.log]

  [Internal]
  _get_api_client()                    -> google.genai.Client  [lazy singleton]
  _call_realtime_api(phase, model, system, user,
                     max_tokens, retries=10) -> str
  # temperature and thinking_budget derived from phase:
  #   phase=="fitness" -> temperature=0.0, thinking_budget=512
  #   all other phases -> temperature=0.7, thinking_budget=2048
  _call(phase, model, system, user, api_max_tokens) -> str  [thin wrapper]
\end{lstlisting}

\section{\fn{\_get\_api\_client()} --- Lazy Singleton}

\begin{spec}{\_get\_api\_client() -> google.genai.Client}
Called by: \_call\_realtime\_api() on first use \\
Thread-safe: double-checked locking with \_client\_lock
\end{spec}

\begin{lstlisting}[style=py]
_api_client = None
_api_client_lock = threading.Lock()

def _get_api_client():
    global _api_client
    if _api_client is not None:
        return _api_client          # fast path: already initialised
    with _api_client_lock:
        if _api_client is not None:
            return _api_client      # double-check after acquiring lock
        sa_info    = json.loads(config.SERVICE_ACCOUNT_PATH.read_text())
        project_id = sa_info["project_id"]
        credentials = service_account.Credentials.from_service_account_info(
            sa_info,
            scopes=["https://www.googleapis.com/auth/cloud-platform"])
        _api_client = genai.Client(
            vertexai=True,
            project=project_id,
            location=config.GCP_REGION,
            credentials=credentials,
            http_options=HttpOptions(timeout=120_000),  # 120 s per request
        )
    return _api_client
\end{lstlisting}

The singleton is initialised once and shared across all threads.  No re-authentication
happens during a run.

\section{\fn{\_call\_realtime\_api()} --- The Core Wrapper}

\begin{spec}{\_call\_realtime\_api(phase, model, system, user, max\_tokens, retries=10) -> Optional[str]}
Calls: \_get\_api\_client(), client.models.generate\_content(), tracker.record(), time.sleep()\\
Retries: 10 attempts, exponential backoff 5s $\times$ 2\textasciicircum{attempt} + jitter\\
Returns: response text; \fn{""} if all retries return empty; \fn{None} if all retries fail
\end{spec}

Temperature and thinking budget are derived internally from the \fn{phase} parameter:
\fn{phase=="fitness"} $\Rightarrow$ temperature=0.0, thinking\_budget=512;
all other phases $\Rightarrow$ temperature=0.7, thinking\_budget=2048.

\begin{lstlisting}[style=py, caption={\_call\_realtime\_api() -- Vertex AI invocation (abridged)}]
def _call_realtime_api(phase, model, system, user,
                       max_tokens, retries=10):
    client = _get_api_client()

    for attempt in range(retries):
        t0 = time.monotonic()
        try:
            is_fitness = (phase == "fitness")
            response = client.models.generate_content(
                model=model,
                contents=user,
                config=GenerateContentConfig(
                    system_instruction=system,
                    max_output_tokens=max_tokens,
                    temperature=0.0 if is_fitness else 0.7,
                    thinking_config=ThinkingConfig(
                        thinking_budget=512 if is_fitness else 2048,
                    ),
                ),
            )
            elapsed = round(time.monotonic() - t0, 2)
            text = response.text or ""

            # token accounting -- thinking tokens billed at output rate
            usage     = response.usage_metadata
            in_toks   = getattr(usage, "prompt_token_count",     0) or 0
            out_toks  = getattr(usage, "candidates_token_count", 0) or 0
            think_toks= getattr(usage, "thoughts_token_count",   0) or 0
            out_toks += think_toks     # merge thinking into output

            cost = _compute_gemini_cost(model, in_toks, out_toks)
            tracker.record(phase=phase, model=model,
                           input_tokens=in_toks, output_tokens=out_toks,
                           cost_usd=cost)

            if not text.strip():            # empty response -> retry
                wait = 5 * (2**attempt) + random.uniform(0, 3)
                time.sleep(wait); continue

            return text                     # SUCCESS

        except Exception as exc:
            # classify: rate-limit (429), server error (500/503), timeout
            # all retryable: wait = 5 * 2^attempt + jitter; time.sleep(wait)
            wait = 5 * (2 ** attempt) + random.uniform(0, 3)
            time.sleep(wait)

    return None                             # all retries exhausted
\end{lstlisting}

\begin{concept}{What happens when \fn{\_call\_realtime\_api} raises \fn{LLMCallError}}
The caller (\fn{generate\_code} etc.) propagates the exception to
\fn{\_eval\_code\_for\_problem()}, which catches it and returns \fn{(0.0, None)}.
The problem scores 0; the run continues.  One failed API call does not crash the GA.
\end{concept}

\section{\fn{generate\_code()} --- Fitness Evaluation}

\begin{spec}{generate\_code(system\_prompt, question) -> str}
Called by: ga.\_eval\_code\_for\_problem(); main.evaluate\_on\_holdout() \\
Calls: \_call\_realtime\_api(phase="fitness", model=TARGET\_MODEL,
thinking\_budget=512, temperature=0.0)
\end{spec}

\begin{lstlisting}[style=py]
_CODE_GEN_DIRECTIVE = (
    "\n\n"
    "ABSOLUTE OUTPUT REQUIREMENT -- NO EXCEPTIONS:\n"
    "Your entire response must consist of ONLY the raw C++ source code -- "
    "one complete file, nothing else. Use C++23.\n"
    "FORBIDDEN:\n"
    "  * Backtick code fences of any form (no ```, no `, no 'cpp' or 'c++' tags)\n"
    "  * Markdown formatting of any kind\n"
    "  * Any text before the code: no explanations, no preamble\n"
    "  * Any text after the code: no summaries, no notes\n"
    "The FIRST character of your response must be the first character of the C++ code.\n"
    "The LAST character must be the last character of the C++ code."
)

def _build_fitness_user_msg(question: str) -> str:
    return (
        f"{question}\n\n"
        "Write a complete, compilable C++ program that reads from stdin and writes "
        "the answer to stdout. Include main(). Compile target: g++ -O2 -std=c++23.\n\n"
        "DO NOT write any explanation, analysis, or reasoning.\n"
        "DO NOT start with 'Looking at', 'I need to', or any English text.\n"
        "Your response must begin with #include or int main -- nothing before it."
    )

def generate_code(system_prompt, question) -> str:
    user_msg = _build_fitness_user_msg(question)
    result = _call(
        phase="fitness",
        model=config.TARGET_MODEL,              # gemini-2.5-flash
        system=system_prompt + _CODE_GEN_DIRECTIVE,
        user=user_msg,
        api_max_tokens=config.API_FALLBACK_MAX_TOKENS_CODE,  # 8192
    )
    if result is None:
        raise LLMCallError(
            f"generate_code: all retries failed (model={config.TARGET_MODEL})")
    return result
\end{lstlisting}

\section{\fn{crossover()} and the Three Mutation Operators}

\begin{lstlisting}[style=chain]
crossover(prompt_a, prompt_b) -> str        [phase="crossover", OPTIMIZER_MODEL, Pro]
  _call(phase="crossover", model=OPTIMIZER_MODEL,
        api_max_tokens=API_FALLBACK_MAX_TOKENS_OPS=700)
  -- temperature=0.7, thinking_budget=2048 derived from phase != "fitness"
  System: "You merge coding-instruction prompts. Output ONLY the merged prompt."
  User:   "Combine ... preserve strongest strategies ...
           --- Parent A --- {prompt_a}
           --- Parent B --- {prompt_b}"

mutate_inject(prompt, failure_examples=None) -> str  [phase="mutate_inject", Pro]
  If failure_examples provided: appends up to 3 formatted failure entries
  (problem_id, rating, reason, compile_error/expected/actual) to guide the Optimizer.
  "Add one new, specific coding strategy ... (address at least one failure pattern)"
  Without failure_examples: generic inject -- "Add one new coding strategy ..."

mutate_delete(prompt) -> str               [phase="mutate_delete", Pro]
  Edge case FIRST: sentences = re.split(r"(?<=[.!?])\s+", prompt)
  if len(sentences) <= 1: return prompt  # single-sentence prompt; skip API call
  "Remove one sentence -- the one that contributes the least to code-generation quality."

mutate_rephrase(prompt) -> str             [phase="mutate_rephrase", Pro]
  "Rephrase using different words, keep same meaning and all same strategies."
  Output: only the rephrased prompt text (no labels, no markdown).
\end{lstlisting}

\begin{concept}{\fn{failure\_examples} in \fn{mutate\_inject}}
Each \fn{Individual} stores up to 5 diverse failure examples
(\fn{individual.failure\_examples}: a list of problem dicts with question, expected
output, actual output).  These are passed to \fn{mutate\_inject()} so the Optimizer
can craft a strategy that specifically addresses the weaknesses the current prompt
exhibits --- targeted mutation rather than random injection.
\end{concept}


% ==========================================================================
\chapter{\pyf{ga.py} --- The Genetic Algorithm}
% ==========================================================================

\section{Logging Infrastructure}

\begin{lstlisting}[style=py]
# Two log streams per run:
#   1. progress.log              -- human-readable timeline (thread-safe)
#   2. evaluations_genNN.jsonl   -- one JSON line per problem evaluation

_progress_log: Path | None = None
_eval_log_dir: Path | None = None
_log_lock      = threading.Lock()
_eval_log_lock = threading.Lock()

def _log(msg, level="INFO"):      # timestamped, level-tagged; prints + appends
def _debug(msg):                   # shortcut for DEBUG-level entries
def _log_evaluation(entry):       # appends one JSON line to evaluations_genNN.jsonl
\end{lstlisting}

\section{\fn{Individual} --- The Chromosome Dataclass}

\begin{lstlisting}[style=py]
@dataclass
class Individual:
    prompt:    str             # the system prompt (the chromosome)
    fitness:   float = -1.0   # -1.0 = sentinel: not yet evaluated
    uid:       str = field(default_factory=lambda: uuid.uuid4().hex[:8])
    parent_a:  str = ""        # uid of first parent  (empty for seeds)
    parent_b:  str = ""        # uid of second parent (empty for seeds)
    operators: str = ""        # e.g. "crossover+inject" or "seed"
    intermediate_steps: dict = field(default_factory=dict)
    # intermediate_steps: per-problem score breakdown populated by compute_fitness
    failure_examples:   list = field(default_factory=list)
    # failure_examples: up to 5 {question, expected, actual} dicts from failed tests
\end{lstlisting}

\section{RNG State Serialisation}

\begin{lstlisting}[style=py]
def rng_state_to_dict(state: tuple) -> dict:
    """Serialise a raw MT19937 state tuple to a JSON-safe dict."""
    version, internalstate, gauss_next = state
    return {
        "version":       version,
        "internalstate": list(internalstate),   # tuple -> list for JSON
        "gauss_next":    gauss_next,
    }

def rng_state_from_dict(d: dict) -> tuple:
    """Restore raw state tuple from a dict saved by rng_state_to_dict."""
    return (d["version"], tuple(d["internalstate"]), d["gauss_next"])

# Usage in run():
#   Save:    checkpoint["rng_state"] = rng_state_to_dict(rng.getstate())
#   Restore: rng.setstate(rng_state_from_dict(checkpoint["rng_state"]))
\end{lstlisting}

These functions enable exact reproduction of any resumed run --- the RNG state is
saved in every \fn{generation\_NN.json} checkpoint so that generation N+1 draws
the same random numbers whether resumed or run fresh.

\section{\fn{\_eval\_one\_problem()} and \fn{\_eval\_code\_for\_problem()} --- Per-Thread Workers}

Evaluation is split across two functions:

\begin{spec}{\_eval\_one\_problem(gen, individual, prob, idx, total) -> (float, dict|None)}
Called by: compute\_fitness() via ThreadPoolExecutor \\
Calls: llm.generate\_code(), then \_eval\_code\_for\_problem()
\end{spec}

\begin{lstlisting}[style=chain]
_eval_one_problem(gen, individual, prob, idx, total):
  time.sleep(random.uniform(0, 1))   # stagger concurrent API requests

  t_llm = time.monotonic()
  try:
      response = llm.generate_code(individual.prompt, prob["question"])
  except LLMCallError as exc:
      llm_time = round(monotonic() - t_llm, 3)
      return _eval_code_for_problem(gen, individual, prob, idx, total,
                                    None, llm_time)   # None = infra failure

  llm_time = round(monotonic() - t_llm, 3)
  return _eval_code_for_problem(gen, individual, prob, idx, total,
                                response, llm_time)
\end{lstlisting}

\begin{spec}{\_eval\_code\_for\_problem(gen, individual, prob, idx, total, response\_text, llm\_time\_s) -> (float, dict|None)}
Called by: \_eval\_one\_problem() \\
Calls: evaluator.extract\_code(), evaluator.run\_tests\_verbose(), \_log\_evaluation()
\end{spec}

\begin{lstlisting}[style=chain]
_eval_code_for_problem(gen, individual, prob, idx, total, response_text, llm_time_s):
  if response_text is None:             # infra failure path
      _log_evaluation({..., score:-1.0, reason:"infra_failure"})
      return (-1.0, None)

  code = evaluator.extract_code(response_text); del response_text

  if not code or not code.strip():
      test_result = {score:0.0, reason:"empty_response", ...}
  else:
      test_result = evaluator.run_tests_verbose(code, prob["inputs"],
                                                prob["outputs"], EXEC_TIMEOUT)

  _log_evaluation({ts, generation, individual_uid, individual_operators,
                   system_prompt, problem_id, problem_rating, llm_call_time_s,
                   extracted_code, compile, test_cases, tests_total, tests_run,
                   early_stopped, total_test_time_s, score, score_reason})
  del entry, code

  # build failure_example if score < 1.0 (and not infra_failure)
  failure_example = {problem_id, rating, reason,
                     expected[:200], actual[:200], compile_error[:200]}

  return (test_result["score"], failure_example)   # (float, dict|None)
\end{lstlisting}

\section{\fn{\_assign\_failure\_examples(ind, collected\_failures)}}

\begin{spec}{\_assign\_failure\_examples(ind, collected\_failures) -> None}
Called by: compute\_fitness() after all threads complete \\
Modifies: ind.failure\_examples (up to 5 diverse examples)
\end{spec}

Selects up to 5 failure examples from \fn{collected\_failures}, prioritising
diversity by failure reason using a priority ordering:

\begin{lstlisting}[style=py]
def _assign_failure_examples(individual, collected_failures):
    _reason_priority = {
        "compile_fail":        0,   # highest priority -- fix compilation first
        "all_fail":            1,
        "all_fail_early_stop": 1,
        "partial_pass":        2,   # lowest priority
    }
    collected_failures.sort(
        key=lambda f: _reason_priority.get(f["reason"], 3))
    seen_reasons: set[str] = set()
    diverse: list[dict] = []
    for f in collected_failures:
        if f["reason"] not in seen_reasons or len(diverse) < 3:
            diverse.append(f)
            seen_reasons.add(f["reason"])
        if len(diverse) >= 5:
            break
    individual.failure_examples = diverse
\end{lstlisting}

Compile failures are prioritised highest --- a prompt that produces uncompilable
code should address that before worrying about partial passes.  After sorting,
each unique failure reason is admitted first; once \fn{diverse} has 3+ entries,
only new reasons are admitted until the cap of 5 is reached.  These are stored
on the \fn{Individual} and passed to \fn{llm.mutate\_inject()} to target real
weaknesses.

\section{\fn{compute\_fitness()} $\to$ \fn{float}}

\begin{spec}{compute\_fitness(gen, individual, problems) -> float in [0.0, 1.0]}
Called by: run() -- 5 times at gen 0; 5 times per generation (gens 1--12)
\end{spec}

\begin{lstlisting}[style=py]
def compute_fitness(gen, individual, problems):
    scores = [0.0] * len(problems)
    collected_failures = []
    _failures_lock = threading.Lock()

    with ThreadPoolExecutor(max_workers=config.PARALLEL_EVALS) as executor:
        futures = {
            executor.submit(
                _eval_one_problem, gen, individual, prob, i, len(problems)
            ): i
            for i, prob in enumerate(problems)
        }
        for future in as_completed(futures):
            idx = futures[future]
            try:
                score, failure_example = future.result()
                scores[idx] = score
                if failure_example is not None:
                    with _failures_lock:
                        collected_failures.append(failure_example)
            except Exception as exc:
                scores[idx] = -1.0

    _assign_failure_examples(individual, collected_failures)
    return sum(scores) / len(scores)
\end{lstlisting}

\begin{warn}{Expensive operation}
One call to \fn{compute\_fitness()} makes \textbf{65 Gemini Flash API calls}.
With \fn{PARALLEL\_EVALS=5}, these run in 5-wide batches.  Evaluating all
$5 \times 13 = 65$ new individuals in a full run (including gen 0 seeds) takes
several hours.
\end{warn}

\section{\fn{tournament\_select}, \fn{roulette\_select}, and \fn{breed}}

\begin{lstlisting}[style=chain]
tournament_select(population, rng):
  contestants = rng.sample(population, TOURNAMENT_SIZE=3)
  return max(contestants, key=lambda ind: ind.fitness)

roulette_select(population, rng):
  fitnesses = [max(ind.fitness, 0.0) for ind in population]
  total = sum(fitnesses)
  if total == 0:                    # all fitnesses zero: fall back to uniform
      return rng.choice(population)
  r = rng.uniform(0, total)
  cumulative = 0.0
  for ind, f in zip(population, fitnesses):
      cumulative += f
      if cumulative >= r:
          return ind
  return population[-1]             # fallback if floating-point rounds short

breed(parent_a, parent_b, rng, allow_mutation=True):
  [BUILD COMBINED FAILURE LIST -- deduplicated union from both parents]
  combined_failures = []
  seen_ids = set()
  for f in (parent_a.failure_examples + parent_b.failure_examples):
      if f["problem_id"] not in seen_ids:
          combined_failures.append(f); seen_ids.add(f["problem_id"])
  combined_failures = combined_failures[:5]

  [CROSSOVER or COPY]
  if rng.random() < 0.80:
      prompt = llm.crossover(parent_a.prompt, parent_b.prompt); ops=["crossover"]
  else:
      prompt = rng.choice([parent_a.prompt, parent_b.prompt]); ops=["copy"]

  [MUTATIONS -- only if allow_mutation is True]
  if allow_mutation:
      if rng.random() < 0.25:
          prompt = llm.mutate_inject(prompt, combined_failures)
          ops += ["inject"]
      if rng.random() < 0.15:
          prompt = llm.mutate_delete(prompt); ops += ["delete"]
      if rng.random() < 0.10:
          prompt = llm.mutate_rephrase(prompt); ops += ["rephrase"]

  return Individual(prompt, parent_a=parent_a.uid, parent_b=parent_b.uid,
                    operators="+".join(ops))   # fitness=-1.0 (unevaluated)
\end{lstlisting}

\section{\fn{run()} --- The GA Loop}

\begin{spec}{run(seed\_prompts, problems, run\_dir, on\_generation,
resume\_population, resume\_start\_gen, resume\_history, resume\_rng\_state,
selection\_scheme="tournament", elitism\_count=None)
-> (best\_Individual, history: list[dict])}
Called by: main.main()
\end{spec}

\begin{lstlisting}[style=chain]
run(seed_prompts, problems, run_dir, on_generation,
    resume_population=None, resume_start_gen=0,
    resume_history=None, resume_rng_state=None,
    selection_scheme="tournament", elitism_count=None):

  _progress_log = run_dir / "progress.log"
  _eval_log_dir = run_dir
  history = resume_history or []

  [RESOLVE SELECTION SCHEME]
  _scheme = selection_scheme.lower()
  _select = roulette_select if _scheme == "roulette" else tournament_select
  _elitism = elitism_count if elitism_count is not None else config.ELITISM_COUNT

  rng = random.Random(config.RANDOM_SEED)   # seeded MT19937

  if resume_population:
      population = resume_population        # restore from checkpoint
      rng.setstate(rng_state_from_dict(resume_rng_state))  # exact MT19937 restore
      loop_start = resume_start_gen + 1
  else:
      population = [Individual(prompt=p, operators="seed")
                    for p in seed_prompts[:POPULATION_SIZE]]
      # Per-seed crash recovery: reload any seed_NN_*.json checkpoints
      # from a previous interrupted gen-0 evaluation
      for i, ind in enumerate(population):
          if <seed checkpoint exists for seed i>:
              restore ind.uid, ind.fitness, ind.failure_examples from ckpt
          else:
              ind.fitness = compute_fitness(0, ind, problems)
              write seed checkpoint: {seed_idx, uid, fitness, failure_examples}
      loop_start = 1

  [GENS loop_start..GENERATIONS]
  stagnation_count = 0
  for gen in range(loop_start, GENERATIONS + 1):
      population.sort(by fitness, descending)
      stats = {generation:gen-1, best, mean, worst, std, cv, best_uid, best_prompt}
      history.append(stats)

      [CONVERGENCE CHECK]
      cv = std / mean  (if mean > CONVERGENCE_MIN_MEAN, else -1.0)
      if cv >= 0 and cv < CONVERGENCE_CV_THRESHOLD:
          stagnation_count += 1
          if stagnation_count >= CONVERGENCE_PATIENCE:
              break   # early stop
      else:
          stagnation_count = 0

      if on_generation:
          on_generation(gen-1, population, stats, rng.getstate())  # saves generation_NN.json

      [ELITISM -- copy top _elitism individuals unchanged]
      new_pop = population[:_elitism]   # typically 0 (tournament/roulette) or 1 (elitism scheme)

      [BREED n_offspring = POPULATION_SIZE - _elitism offspring]
      offspring_idx = 0
      while len(new_pop) < POPULATION_SIZE:
          p1 = _select(population, rng)   # tournament or roulette per --scheme
          p2 = _select(population, rng)
          allow_mut = (offspring_idx >= CROSSOVER_ONLY_SLOTS)  # first 2: crossover only
          child = breed(p1, p2, rng, allow_mutation=allow_mut)
          new_pop.append(child); offspring_idx += 1

      [EVALUATE offspring only (elites already have valid fitness)]
      for ind in new_pop[_elitism:]:
          ind.fitness = compute_fitness(gen, ind, problems)

      population = new_pop

  [FINAL] population.sort()
  history.append({generation:last_completed_gen, ...})
  return (population[0], history)
\end{lstlisting}

\begin{concept}{Three selection schemes (\fn{--scheme} argument)}
\fn{tournament} (default): sample 3 individuals at random, pick the fittest.
Balances exploitation with exploration; independent of absolute fitness values.\\
\fn{roulette}: fitness-proportionate selection; probability of picking an individual
is proportional to its fitness.  Falls back to uniform random when all fitnesses
are zero.\\
\fn{elitism}: tournament selection + 1 elite individual copied unchanged each
generation (\fn{elitism\_count=1}).  Guarantees the best-ever individual survives.
\end{concept}

\begin{concept}{Convergence detection (CV)}
The coefficient of variation $CV = \sigma / \mu$ measures population homogeneity.
When $CV < 0.01$ (less than 1\% variation) for 3 consecutive generations and
$\mu > 0.01$ (not a trivial all-fail population), the GA stops early.
This prevents wasting API budget on a converged population.
\end{concept}

\begin{concept}{Per-seed checkpoints for crash recovery}
Each seed's fitness is saved to \fn{seed\_NN\_\textless uid\textgreater.json} before
moving to the next seed.  If the run crashes midway through gen-0 evaluation, it
re-loads completed seeds from disk rather than re-evaluating them.
\end{concept}


% ==========================================================================
\chapter{\pyf{main.py} --- Orchestration \& Entry Point}
% ==========================================================================

\section{The Key Functions}

\subsection{Vertex AI Preflight (replaces \fn{\_check\_cli()})}

\begin{spec}{Vertex AI credential check in main()}
Called at startup, before dataset loading \\
Calls: Path.exists(), \_get\_api\_client() (via llm module import)
\end{spec}

Instead of checking for the \fn{claude} CLI executable, \fn{main()} verifies that
\fn{config.SERVICE\_ACCOUNT\_PATH} (\fn{service-account.json}) exists on disk.  It
then imports \fn{llm}, which triggers \fn{\_get\_api\_client()} --- if the service
account is invalid or the project lacks Vertex AI permissions, the error surfaces
immediately before any dataset work.

\subsection{\fn{\_find\_latest\_run(results\_dir)} and \fn{\_load\_resume\_state()}}

\begin{spec}{\_find\_latest\_run(results\_dir) -> Path | None}
Called by: main() when \fn{--resume} flag is given without an explicit path \\
Returns: the most recently modified run directory under \fn{results/}
\end{spec}

\begin{spec}{\_load\_resume\_state(run\_dir) -> dict}
Called by: main() after finding the run to resume \\
Returns: \{population, start\_gen, history, rng\_state\} loaded from the latest
generation\_NN.json checkpoint
\end{spec}

Resume support allows a crashed run to continue from the last saved generation,
replaying \fn{calls.jsonl} into \fn{tracker.restore\_from\_jsonl()} to restore
accurate running cost totals.

\subsection{\fn{save\_checkpoint(run\_dir, gen, population, stats, rng)}}

\begin{spec}{save\_checkpoint(...)}
Called by: the on\_generation lambda in main() \\
Output: results/run\_*/generation\_NN.json
\end{spec}

Saves the full population state (all 5 prompts, fitnesses, parent UIDs, operators,
\fn{failure\_examples}, \fn{intermediate\_steps}) plus the RNG state after each
generation.  The RNG state (MT19937 internals) is serialised via
\fn{rng\_state\_to\_dict()} so that a resumed run produces identical subsequent draws.

\subsection{\fn{evaluate\_on\_holdout(best\_prompt, holdout)} $\to$ \fn{float}}

\begin{spec}{evaluate\_on\_holdout(best\_prompt, holdout, run\_dir) -> float}
Called by: main() \\
Calls: ThreadPoolExecutor, llm.generate\_code(), evaluator.extract\_code(),
evaluator.run\_tests\_verbose(); writes holdout\_evaluations.jsonl
\end{spec}

Now runs \textbf{in parallel} using \fn{ThreadPoolExecutor(max\_workers=PARALLEL\_EVALS)}
(was sequential in the previous version).  Evaluates all 200 held-out problems with
the best evolved prompt; returns the mean partial-credit score.

\subsection{\fn{main()} --- The 8-Step Pipeline}

\begin{lstlisting}[style=chain]
main():

  [0] Parse CLI args
      --resume [RUN_DIR]          (optional: resume from crash)
      --scheme {tournament|roulette|elitism}  (default: "tournament")
      --help

  [1] Pre-flight checks
      verify SERVICE_ACCOUNT_PATH exists
      verify APPS_TRAIN / APPS_TEST paths exist (or warn)
      verify Vertex AI connectivity (lazy via llm module init)

  [2] Load and split dataset
      all_problems = dataset.load_problems()          # 305 Codeforces problems
      evolution_pool, holdout, _ = dataset.split(
          all_problems, evolution_size=65,
          holdout_size=200, seed=42)

  [3] Create timestamped run directory (or locate existing for resume)
      run_dir = RESULTS_DIR / "run_YYYYMMDD_HHMMSS_<scheme>"
                e.g. "run_20240501_143022_tournament"
      Writes: config.json           (hyperparameter snapshot)
              evolution_ids.json    (which 65 problem IDs)

  [4] Initialise logging and cost tracking
      tracker.set_log_path(run_dir / "calls.jsonl")
      tracker.set_live_path(run_dir / "costs_live.json")
      tracker.set_costs_live_path(COSTS_DIR / f"cost_{run_dir.name}.json")
      llm.set_error_log_path(run_dir / "llm_errors.log")
      if resuming: tracker.restore_from_jsonl(run_dir / "calls.jsonl")

  [5] Run GA
      resume_state = _load_resume_state(run_dir) if resuming else {}
      best, history = ga.run(
          seed_prompts=SEED_PROMPTS[:POPULATION_SIZE],   # first 5
          problems=evolution_pool,
          run_dir=run_dir,
          on_generation=_on_generation,    # saves generation_NN.json + logs cost
          resume_population= resume_state["population"]  if resume_state else None,
          resume_start_gen=  resume_state["start_gen"]   if resume_state else 0,
          resume_history=    resume_state["history"]     if resume_state else None,
          resume_rng_state=  resume_state["rng_state"]   if resume_state else None,
          selection_scheme=args.scheme,
          elitism_count=1 if args.scheme == "elitism" else None)

  [6] Save results
      (run_dir / "history.json").write_text(json.dumps(history))
      (run_dir / "best_prompt.txt").write_text(best.prompt)

  [7] Holdout evaluation (parallel)
      holdout_fitness = evaluate_on_holdout(best.prompt, holdout)
      (run_dir / "summary.json").write_text({
          elapsed_minutes, best_prompt, best_fitness_evolution,
          best_fitness_holdout, best_uid, total_generations,
          resumed, selection_scheme})

  [8] Genealogy report + cost flush
      genealogy.generate_report(run_dir, best.uid)
      -> run_dir/genealogy_report.txt, run_dir/genealogy_tree.json
      tracker.flush(extra={...}) -> Code/Costs/cost_TIMESTAMP.json
\end{lstlisting}

\section{Output Files}

\begin{center}
\begin{tabular}{@{}p{5.5cm}p{6.7cm}@{}}
\toprule
\textbf{File} & \textbf{Contents} \\
\midrule
\fn{run\_*/config.json}              & Exact hyperparameters (reproducibility) \\
\fn{run\_*/evolution\_ids.json}      & The 65 problem IDs used as the fitness pool \\
\fn{run\_*/calls.jsonl}              & Per-call token + cost log (streaming JSONL) \\
\fn{run\_*/costs\_live.json}         & Live running-total cost summary \\
\fn{run\_*/llm\_errors.log}          & API retry / error events \\
\fn{run\_*/progress.log}             & Human-readable timestamped timeline \\
\fn{run\_*/evaluations\_genNN.jsonl} & Per-problem full evaluation records \\
\fn{run\_*/seed\_NN\_\textless uid\textgreater.json} & Per-seed gen-0 checkpoint (crash recovery) \\
\fn{run\_*/generation\_NN.json}      & Full population + RNG state per generation \\
\fn{run\_*/history.json}             & Best / mean / worst fitness per generation \\
\fn{run\_*/best\_prompt.txt}         & The evolved prompt (main artefact) \\
\fn{run\_*/holdout\_evaluations.jsonl} & Per-problem results for the 200 holdout problems \\
\fn{run\_*/summary.json}             & Final fitnesses (evolution + holdout), scheme, runtime \\
\fn{run\_*/genealogy\_report.txt}    & Human-readable prompt ancestry chain \\
\fn{run\_*/genealogy\_tree.json}     & Full ancestry tree in JSON \\
\fn{Costs/cost\_*.json}              & Final token counts and estimated cost breakdown per run \\
\fn{Costs/cost\_\textless run\_name\textgreater.json} & Live cost summary updated after every API call (named after run dir) \\
\bottomrule
\end{tabular}
\end{center}


% ==========================================================================
\chapter{\pyf{genealogy.py} --- Ancestry Tracing}
% ==========================================================================

\section{Purpose}

After a GA run completes, \pyf{genealogy.py} traces the lineage of the best evolved
prompt back through all checkpoints to its seed ancestor.  The output is both a
human-readable report and a JSON tree.

\section{The Three Core Functions}

\begin{lstlisting}[style=chain]
_load_all_individuals(run_dir) -> dict[uid, dict]
  Scans run_dir for all generation_NN.json checkpoint files.
  Builds a registry: {uid: {prompt, fitness, parent_a, parent_b,
                            operators, generation}}.
  Later generations overwrite earlier entries for the same uid (fitness may
  have been updated by re-evaluation on resume).

_collect_lineage(uid, registry) -> list[dict]
  Walks the parent_a chain from the best individual back to a seed.
  A seed has parent_a == "" (empty string).
  Returns a list of individual dicts ordered from seed to best.

_collect_full_tree(uid, registry, depth=0) -> dict
  Recursively collects both parents (parent_a and parent_b) at each node.
  Depth-limited to 40 to prevent infinite loops from any cycle in the graph.
  Returns a nested dict suitable for JSON serialisation.

generate_report(run_dir, best_uid) -> None
  Entry point called by main.main().
  1. _load_all_individuals(run_dir)
  2. _collect_lineage(best_uid, registry)    -> linear chain
  3. _collect_full_tree(best_uid, registry)  -> full ancestry tree
  4. Writes run_dir/genealogy_report.txt  (human-readable text)
  5. Writes run_dir/genealogy_tree.json   (full tree dict)
\end{lstlisting}

\section{Genealogy Report Format}

\begin{lstlisting}[style=chain]
=== Genealogy Report ===
Best individual: uid=a3f2c1b9  fitness=0.7231  generation=11

Linear ancestry (seed -> best):
  Gen 0  uid=seed_1   fitness=0.4102  operators=seed
         prompt: "Solve the following C++ competitive programming..."
  Gen 3  uid=b2d44f1a fitness=0.5810  operators=crossover+inject
         parents: [seed_1 x seed_3]
         prompt: "Analyse the problem carefully. Consider time..."
  ...
  Gen 11 uid=a3f2c1b9 fitness=0.7231  operators=crossover+inject
         parents: [e1f2a3b4 x c5d6e7f8]
         prompt: "You are a competitive programmer. First identify..."

Full ancestry tree written to: run_dir/genealogy_tree.json
\end{lstlisting}

\begin{concept}{Why genealogy matters}
The linear ancestry chain shows which genetic operations were most productive.
If the best prompt descends through 3 consecutive \fn{crossover+inject} operations,
that suggests inject mutations are the key driver.  If the chain is mostly
\fn{crossover+copy}, injection may not be contributing.  This guides future
hyperparameter choices.
\end{concept}


% ==========================================================================
\chapter*{Glossary}
\addcontentsline{toc}{chapter}{Glossary}
% ==========================================================================

\begin{description}[style=nextline, leftmargin=3cm, itemsep=2pt]

\item[Chromosome] A candidate solution in a GA.  Here: a system prompt string,
  stored as an \fn{Individual} object.

\item[Convergence detection] Stopping the GA early when the population is
  statistically homogeneous: $CV = \sigma/\mu < 0.01$ for 3 consecutive generations
  and $\mu > 0.01$.  Prevents wasting API budget on a converged population.

\item[CROSSOVER\_ONLY\_SLOTS] The first $N=2$ offspring bred per generation receive
  crossover only (no mutation), providing stability without hard-locking elite
  individuals.  The remaining 3 offspring per generation are mutation-eligible.

\item[Crossover] Genetic operator combining two parents into one offspring.  Here:
  the Optimizer LLM semantically blends two parent prompts.  Fires with probability 0.8.

\item[CV (coefficient of variation)] $\sigma / \mu$ of fitness across the population.
  Used by convergence detection.

\item[Early stopping (test loop)] In \fn{run\_tests\_verbose()}: halting the test loop
  after \fn{EARLY\_STOP\_FAILURES=3} consecutive failures.  The score denominator is
  still the total test count, so skipped tests count as failures.

\item[Elitism] Copying the top $k$ individuals unchanged into the next generation.
  \fn{ELITISM\_COUNT=0} in the current version --- no frozen elites.

\item[Evolution pool] The fixed 65 Codeforces problems used as the fitness benchmark.
  Chosen by stratified shuffle with seed 42 (5 harder problems always included).

\item[Failure examples] Up to 5 problem failures stored on each \fn{Individual}
  (\fn{individual.failure\_examples}).  Passed to \fn{mutate\_inject()} so the
  Optimizer crafts strategies that address real weaknesses.

\item[Fitness] Mean partial-credit score across 65 problems.  Range: [0.0, 1.0].
  Computed by \fn{ga.compute\_fitness()}.

\item[Genealogy] The ancestry chain from the best evolved prompt back to its seed.
  Traced by \fn{genealogy.py} after a run completes.

\item[Gemini Flash] \fn{gemini-2.5-flash} --- the Target model used for code
  generation (fitness evaluation).  Fast, cheap.

\item[Gemini Pro] \fn{gemini-2.5-pro} --- the Optimizer model used for crossover
  and mutation meta-prompts.  More capable, used for creative prompt manipulation.

\item[Held-out set] 200 problems never seen during evolution.  Used to test
  generalisation of the final evolved prompt.

\item[Individual] One GA population member: a prompt with fitness, uid, parent uids,
  operator history, intermediate steps, and failure examples.

\item[Intermediate steps] Per-problem scores stored in \fn{individual.intermediate\_steps}
  as a dict mapping problem index to score.  Used for analysis and checkpoint saves.

\item[Meta-prompt] A prompt sent to the Optimizer LLM instructing it to modify another
  prompt (used in crossover and all three mutation operators).

\item[Mutation] Modifying a single chromosome.  Three independent types: inject (25\%),
  delete (15\%), rephrase (10\%).  Only applies to offspring beyond CROSSOVER\_ONLY\_SLOTS.

\item[PARALLEL\_EVALS] Number of problems evaluated simultaneously per individual
  (\fn{PARALLEL\_EVALS=5}, using a \fn{ThreadPoolExecutor}).

\item[Partial credit] Score = (tests passed) / (total tests).  Provides a smooth
  gradient vs.\ binary pass/fail.

\item[Process group] A set of processes sharing a PGID.  \fn{evaluator.py} creates
  new process groups via \fn{preexec\_fn=os.setsid} and kills them with
  \fn{os.killpg(pgid, SIGKILL)} to prevent orphaned children after timeouts.

\item[Resume] The \fn{--resume} CLI flag re-loads the latest \fn{generation\_NN.json}
  checkpoint and continues the GA from the next generation, replaying cost totals from
  \fn{calls.jsonl}.

\item[RNG state] The internal state of Python's MT19937 Mersenne Twister random number
  generator.  Serialised to JSON via \fn{rng\_state\_to\_dict()} and saved in every
  \fn{generation\_NN.json} checkpoint, enabling exact reproduction of resumed runs.

\item[Service account] A GCP JSON credential file (\fn{service-account.json}) used to
  authenticate Vertex AI API calls.  No interactive login or API key required.

\item[Singleton] The \fn{tracker = CostTracker()} instance at module level is shared
  by all importers (\pyf{llm.py} and \pyf{main.py}).

\item[Startup jitter] A random 0--2 s sleep at the start of each
  \fn{\_eval\_code\_for\_problem()} call to stagger parallel API launches.

\item[Stratified split] Sampling a fixed count from each difficulty tier (rating) when
  building the evolution pool, ensuring representation of all difficulty levels.

\item[System prompt] English instructions sent before the user message, shaping LLM
  behaviour.  The GA chromosome.

\item[Temperature] LLM randomness parameter.  Now enforced via
  \fn{GenerateContentConfig(temperature=...)}: 0.0 for Flash (deterministic code gen),
  0.7 for Pro (creative operator mutations).

\item[Thinking tokens] Internal chain-of-thought tokens produced by Gemini 2.5 models
  before the visible response.  Billed at the output token rate.  Tracked via
  \fn{usage\_metadata.thoughts\_token\_count}.  Budget: 512 for Flash, 2048 for Pro.

\item[Tournament selection] Sample $k=3$ individuals at random; return the fittest.
  Balances exploitation (picking the best) with exploration (not always the global best).

\item[Two-model architecture] Gemini 2.5 Flash (code generation, target) is optimised;
  Gemini 2.5 Pro (reasoning, creative) performs the genetic operators.

\item[Vertex AI] Google Cloud's managed ML platform.  The project uses it to call
  Gemini models via the \fn{google-genai} Python SDK with service account auth.

\end{description}

\end{document}
